\begin{abstract}
Instruction-guided CAD editing must alter the requested feature while preserving the rest of a three-dimensional shape. Program validity and global volume overlap alone cannot establish this behavior. We study executable shape editing with \ours, a code-conditioned pipeline that predicts validated program patches, executes native CAD representations, and evaluates resulting solids and derived views. A fresh 3B model is fine-tuned on 2,540 public edit pairs from BenchCAD and CAD-Editor. On a locally defined 252-task holdout, strict volume agreement increases from 1 to 39 cases; 139 trained outputs execute successfully, exposing a substantial gap between valid geometry and target agreement. On the same holdout, zero-shot GPT-6 Astra reaches \paAstra/252 strict matches. We wrap Astra in an experience-guided graph search with diverse edit hypotheses, a hierarchical geometry, constraint, semantic, FEM and robustness verifier, diagnosis-driven repair, MCTS and Pareto selection. On \pipeN{} BenchCAD tasks it does not raise strict agreement (\pipeAstra{} to \pipeFull), because \pipeWrongPass{} of Astra's \pipeWrong{} incorrect answers pass every target-free check; retrieved training edits alone give \pipeRag, a post hoc, nonsignificant gain. A focused audit of previously illustrated outputs compares four common-camera projections and source-relative edit masks. A flange example demonstrates how high overall silhouette agreement can conceal misplaced mounting holes. A separate diagnostic FEM audit checks four completed prediction/target pairs under explicit synthetic material and load contracts; two other attempts remain unverified. The study contributes an auditable evaluation workflow and empirical evidence, rather than a new CAD architecture. The model receives programs and text, with no learned image input or solver supervision. Local dataset repartitioning, limited visual coverage and a single training seed bound the conclusions.
\end{abstract}
